\documentclass[11pt]{article}

\usepackage[a4paper,margin=1in]{geometry}
\usepackage[T1]{fontenc}
\usepackage{lmodern}
\usepackage{amsmath,amssymb,amsthm}
\usepackage{microtype}
\usepackage{xurl}
\usepackage[hidelinks]{hyperref}

\hypersetup{
  pdftitle={An Accuracy Obstruction for lp Row Sparsification},
  pdfsubject={A complete-graph counterexample to RA-06 for p greater than two},
  pdfkeywords={sensitivity sampling, subspace embedding, row sparsification, graph energy}
}

\linespread{1.06}
\setlength{\parskip}{3pt plus 1pt minus 1pt}
\setlength{\emergencystretch}{2em}
\raggedbottom

\newtheoremstyle{spaced}
  {9pt plus 2pt minus 2pt}{9pt plus 2pt minus 2pt}
  {\itshape}{}{\bfseries}{.}{.5em}{}
\theoremstyle{spaced}
\newtheorem{theorem}{Theorem}[section]
\newtheorem{lemma}[theorem]{Lemma}
\newtheorem{corollary}[theorem]{Corollary}
\numberwithin{equation}{section}

\newcommand{\E}{\mathbb{E}}
\newcommand{\R}{\mathbb{R}}
\newcommand{\Sp}{\mathcal{S}}
\newcommand{\one}{\mathbf{1}}
\DeclareMathOperator{\supp}{supp}

\title{An Accuracy Obstruction for $\ell_p$ Row Sparsification\\[4pt]
  \large A Complete-Graph Counterexample to RA-06}
\author{}
\date{30 September 2026}

\begin{document}
\maketitle

\begin{abstract}
For every fixed real $p>2$, we construct full-column-rank matrices with
ordinary $\ell_p$ sensitivity $\Sp=\Theta_p(d)$ for which every
nonnegative reweighted row subset preserving all $p$-th powers of norms
within $1\pm\varepsilon$ needs
$\Omega_p((\Sp+d)\varepsilon^{-p})$ rows, in an explicit large-dimension
regime. The matrices are grounded incidence matrices of complete graphs.
An elementary, support-dependent test potential proves the lower bound.
For the prescribed ordinary-sensitivity sampler, even one successful
realization forces the same order of expected row count. This disproves
RA-06 and the $p>2$ quadratic-accuracy sensitivity-sampling conjecture it
formalizes: no fixed polylogarithmic factor can absorb the gap between
$\varepsilon^{-p}$ and $\varepsilon^{-2}$.
\end{abstract}

\section{Introduction and the target assertion}\label{sec:introduction}

For a full-column-rank matrix $A\in\R^{n\times d}$ with rows $a_i^T$,
its ordinary $\ell_p$ row sensitivities and total sensitivity are
\begin{equation}\label{eq:sensitivity-definition}
 s_i=\sup_{x\ne0}\frac{|a_i^Tx|^p}{\|Ax\|_p^p},
 \qquad
 \Sp=\sum_{i=1}^n s_i.
\end{equation}
The question is whether these quantities alone support a row-sampling
bound that is nearly linear in $\Sp+d$ and quadratic in the inverse
accuracy.

Woodruff and Yasuda conjectured the sample complexity
$\widetilde O(\varepsilon^{-2}(\Sp+d))$ for sensitivity sampling
\cite[Section~3]{wy}. Their sampling model and ordinary-sensitivity rule
appear in Definition~1.1 and Theorem~1.5 of that work. The problem
catalogue entry RA-06 records the following precise $p>2$ formulation
\cite{ra06}.

\medskip
\noindent\textbf{Target assertion (RA-06).}
Fix $p>2$. Do constants $C_p,c_p>0$ exist such that, for every
full-column-rank $A\in\R^{n\times d}$ and every
$0<\varepsilon,\delta<1/2$, one can choose $\alpha>0$ so that independent
row retention with probabilities
\begin{equation}\label{eq:ra06-probabilities}
 q_i=\min\{1,1/n+s_i/\alpha\}
\end{equation}
satisfies the expected-size bound
\begin{equation}\label{eq:ra06-budget}
 \sum_{i=1}^n q_i
 \le C_p\varepsilon^{-2}(\Sp+d)
      \log^{c_p}\!\left(\frac{2nd}{\varepsilon\delta}\right),
\end{equation}
while the matrix $\widetilde A$ obtained by scaling each retained row by
$q_i^{-1/p}$ obeys, with probability at least $1-\delta$,
\begin{equation}\label{eq:matrix-embedding}
 (1-\varepsilon)\|Ax\|_p^p
 \le \|\widetilde Ax\|_p^p
 \le (1+\varepsilon)\|Ax\|_p^p
 \qquad\text{for every }x\in\R^d?
\end{equation}
\medskip

We give a negative answer for every fixed real $p>2$. The obstruction is
an accuracy cost, not a claim that nearly linear sensitivity dependence
is impossible at fixed accuracy. The argument has two levels. First,
any nonnegative weighted subgraph approximating the complete graph's
$p$-energy must have many support edges. Second, equal sensitivities
force the rule \eqref{eq:ra06-probabilities} to use a common probability,
which converts the deterministic obstruction into a lower bound on the
sampler's expected size. The latter conclusion rules out every allowed
choice of $\alpha$, including the floor and saturation in the rule.

All logarithms are natural. Constants with a subscript $p$ may depend
on the fixed exponent $p$, but not on the matrix dimensions or accuracy.

\section{Complete-graph matrices and exact sensitivities}\label{sec:matrix}

Let $v\ge3$, and write
\[
 [v]=\{1,\ldots,v\},\qquad R=v-1,\qquad
 n=\binom v2,\qquad d=v-1.
\]
Choose an orientation of every edge of the complete graph $K_v$.
Let $A_v\in\R^{n\times d}$ be its edge--vertex incidence matrix with the
column for vertex $v$ deleted. Define the complete-graph $p$-energy by
\begin{equation}\label{eq:complete-energy}
 F(z)=\sum_{1\le i<j\le v}|z_i-z_j|^p,
 \qquad z\in\R^v.
\end{equation}

For $x\in\R^{v-1}$, the grounded potential $z=(x,0)$ satisfies
$\|A_vx\|_p^p=F(z)$. Conversely, replacing any $z\in\R^v$ by
$z-z_v\one$ grounds its last coordinate and leaves every edge difference
unchanged. Thus preservation of $F$ for all vertex potentials is
exactly the matrix embedding requirement for $A_v$; grounding discards
no test direction. The matrix has full column rank, since a grounded
potential whose pairwise differences all vanish is zero.

\begin{lemma}[Exact sensitivities]\label{lem:sensitivity}
For every real $p>1$, every row of $A_v$ has sensitivity
\begin{equation}\label{eq:exact-sensitivity}
 s_v=\frac{1}{1+(v-2)2^{1-p}}.
\end{equation}
Consequently,
\begin{equation}\label{eq:total-sensitivity}
 \Sp=\frac{\binom v2}{1+(v-2)2^{1-p}}
 \le 2^{p-2}v.
\end{equation}
For $p\ge2$, one also has
\begin{equation}\label{eq:sensitivity-comparison}
 d=v-1\le\Sp,
 \qquad
 \Sp+d\le(2^{p-2}+1)v.
\end{equation}
In particular, $\Sp=\Theta_p(d)$ for each fixed $p>2$.
\end{lemma}

\begin{proof}
Fix an edge $\{i,j\}$. For each $w\notin\{i,j\}$, convexity gives
\[
 |z_i-z_w|^p+|z_j-z_w|^p
 \ge 2^{1-p}|z_i-z_j|^p.
\]
Adding these inequalities, the edge $\{i,j\}$, and the remaining
nonnegative energies yields
\[
 F(z)\ge\bigl(1+(v-2)2^{1-p}\bigr)|z_i-z_j|^p.
\]
Equality holds at $z_i=1$, $z_j=0$, and $z_w=1/2$ for every other
vertex. Translating this potential to grounded coordinates proves
\eqref{eq:exact-sensitivity}.

Summing the equal sensitivities gives the identity in
\eqref{eq:total-sensitivity}. Since
$1+(v-2)2^{1-p}\ge(v-1)2^{1-p}$ for $p>1$, its upper bound follows.
For $p\ge2$, the reverse comparison
$1+(v-2)2^{1-p}\le v/2$ gives $\Sp\ge v-1$.
Finally, $d\le v$ gives the last bound in
\eqref{eq:sensitivity-comparison}.
\end{proof}

\section{A lower bound for arbitrary nonnegative row weights}
\label{sec:weighted}

Assign a finite nonnegative weight $w_{ij}$ to each unordered edge
$\{i,j\}$ of $K_v$, and set
\begin{equation}\label{eq:weighted-energy}
 U_w(z)=\sum_{i<j}w_{ij}|z_i-z_j|^p.
\end{equation}
Write $w_{ji}=w_{ij}$ and $w_{ii}=0$. Let $H_w$ be the support graph,
whose edges are those with positive weight. Its number of edges and the
support degree of a vertex $u$ are
\[
 M=|E(H_w)|,
 \qquad
 \kappa_u=|\{j\ne u:w_{uj}>0\}|.
\]
Support degree counts retained edges; it is distinct from the weighted
degree $d_w(u)=\sum_{j\ne u}w_{uj}$.

\begin{theorem}[Deterministic support obstruction]\label{thm:weighted}
Let $p>1$, $v\ge3$, and $0<\varepsilon<1/2$. Suppose that
\begin{equation}\label{eq:weighted-embedding}
 (1-\varepsilon)F(z)\le U_w(z)\le(1+\varepsilon)F(z)
 \qquad\text{for all }z\in\R^v.
\end{equation}
Then every vertex has positive support degree and satisfies
\begin{equation}\label{eq:support-necessary}
 \kappa_u^{-1/p}\le4\varepsilon+\frac{2\kappa_u}{R}.
\end{equation}
In particular,
\begin{equation}\label{eq:support-degree}
 \kappa_u\ge
 \min\{R\varepsilon,(6\varepsilon)^{-p}\},
\end{equation}
and therefore
\begin{equation}\label{eq:weighted-support-bound}
 M\ge\frac v2\min\{R\varepsilon,(6\varepsilon)^{-p}\}.
\end{equation}
\end{theorem}

\begin{proof}
\emph{Singleton tests control weighted degrees.}
Apply \eqref{eq:weighted-embedding} to the potential equal to one at
vertex $u$ and zero elsewhere. This gives
\begin{equation}\label{eq:weighted-degrees}
 (1-\varepsilon)R\le d_w(u)\le(1+\varepsilon)R
 \qquad(u\in[v]).
\end{equation}
The lower bound is positive, so $\kappa_u\ge1$.

\emph{A support-dependent potential exposes the missing edges.}
Fix $u$, let $W=\{j:w_{uj}>0\}$, and write $k=|W|$ and
$O=[v]\setminus(W\cup\{u\})$. Choose
\[
 t=k^{-1/p}\in(0,1],
 \qquad
 z_u=1,\quad z_j=t\ (j\in W),\quad z_j=0\ (j\in O).
\]
The test is permitted to depend on $H_w$: the hypothesis requires
preservation for every potential after the weights are fixed.

There are $k$ edges from $u$ to $W$, $R-k$ from $u$ to $O$, and
$k(R-k)$ from $W$ to $O$. All other complete-graph edges have zero
energy at this potential. Since $kt^p=1$,
\begin{align}
 F(z)
 &=k(1-t)^p+(R-k)+k(R-k)t^p \notag\\
 &=2(R-k)+k(1-t)^p.
 \label{eq:witness-complete}
\end{align}
There are no positive-weight edges from $u$ to $O$. Consequently,
\begin{equation}\label{eq:witness-weighted}
 U_w(z)=d_w(u)(1-t)^p+
       t^p\sum_{i\in W}\sum_{j\in O}w_{ij}.
\end{equation}
The crossing weight is at most
\[
 \sum_{i\in W}\sum_{j\in O}w_{ij}
 \le\sum_{i\in W}d_w(i)
 \le k(1+\varepsilon)R.
\]
Combining this estimate, $t^p=1/k$, and
\eqref{eq:weighted-degrees} at $u$ yields
\begin{equation}\label{eq:witness-upper}
 U_w(z)\le(1+\varepsilon)R\bigl(1+(1-t)^p\bigr).
\end{equation}

\emph{The energy comparison forces a large support degree.}
Put $h=1-(1-t)^p$. The lower half of
\eqref{eq:weighted-embedding}, together with
\eqref{eq:witness-complete} and \eqref{eq:witness-upper}, gives
\[
 2(1-\varepsilon)(R-k)
 \le(1+\varepsilon)R(2-h).
\]
Rearranging,
\begin{equation}\label{eq:nonlinear-necessary}
 (1+\varepsilon)h
 \le4\varepsilon+2(1-\varepsilon)k/R.
\end{equation}
For $p>1$ and $0\le t\le1$, one has $(1-t)^p\le1-t$, hence
$h\ge t$. It follows that
\[
 k^{-1/p}=t\le4\varepsilon+2k/R,
\]
which proves \eqref{eq:support-necessary}.

If $k$ were smaller than both $R\varepsilon$ and
$(6\varepsilon)^{-p}$, the right-hand side would be less than
$6\varepsilon$, while the left-hand side would exceed
$6\varepsilon$. This contradiction proves
\eqref{eq:support-degree}. Summing that inequality over all vertices
and using $\sum_u\kappa_u=2M$ proves
\eqref{eq:weighted-support-bound}.
\end{proof}

The witness lifts every retained neighbor of $u$ to the small common
potential $t$. This reduces all supported edges incident to $u$, whereas
each missing edge from $u$ to $O$ still has difference one in the original
complete graph. Weighted-degree control limits the compensation that
edges leaving $W$ can provide. The balance $kt^p=1$ makes the resulting
error visible on the scale $k^{-1/p}$.

\begin{corollary}[Row-subset and randomized lower bounds]
\label{cor:weighted-accuracy}
Fix $p>2$ and $0<\varepsilon<1/2$. If
\begin{equation}\label{eq:large-v}
 (v-1)\varepsilon\ge(6\varepsilon)^{-p},
 \quad\text{equivalently}\quad
 v-1\ge6^{-p}\varepsilon^{-(p+1)},
\end{equation}
then every nonnegative weighted row subset of $A_v$ satisfying
\eqref{eq:matrix-embedding} has support size
\begin{equation}\label{eq:weighted-accuracy}
 M\ge\frac{6^{-p}}2\,v\varepsilon^{-p}
 \ge\beta_p(\Sp+d)\varepsilon^{-p},
 \qquad
 \beta_p=\frac{6^{-p}}{2(2^{p-2}+1)}.
\end{equation}
If a randomized nonnegative weighted row-subset construction succeeds
with probability at least $1-\delta$, where $0\le\delta<1$, then
\begin{equation}\label{eq:weighted-expectation}
 \E M\ge(1-\delta)\frac{6^{-p}}2\,v\varepsilon^{-p}.
\end{equation}
No independence or restriction on how the weights are chosen is required.
\end{corollary}

\begin{proof}
A weight $w_{ij}$ corresponds to retaining the associated row of $A_v$
scaled by $w_{ij}^{1/p}$. Grounding identifies its matrix energy with
\eqref{eq:weighted-energy}. Apply Theorem~\ref{thm:weighted}, the
condition \eqref{eq:large-v}, and Lemma~\ref{lem:sensitivity} to obtain
\eqref{eq:weighted-accuracy}. Every successful randomized output obeys
the deterministic lower bound, and support size is nonnegative on the
failure event; taking expectations proves
\eqref{eq:weighted-expectation}.
Repeated copies of a row can be combined into one nonnegative weight,
so the lower bounds also apply to the total number of output rows.
\end{proof}

\section{Common-probability sampling: an expected-size obstruction}
\label{sec:common}

For $q\in(0,1]$ and a subgraph $G\subseteq K_v$, define
\[
 U_{G,q}(z)=\frac1q\sum_{\{i,j\}\in E(G)}|z_i-z_j|^p,
 \qquad D=qR.
\]
This is the weighted energy with weight $1/q$ on each retained edge.
If every edge has marginal retention probability $q$, the expected
number of retained edges is exactly
\begin{equation}\label{eq:expected-size}
 \E|E(G)|=nq=\frac{vD}{2}.
\end{equation}
Here $D$ is a deterministic parameter, not a realized graph degree.

\begin{theorem}[Necessary common sampling probability]\label{thm:common}
Let $p>1$, $v\ge3$, $0<\varepsilon<1/2$, and $q\in(0,1]$.
If a subgraph $G$ satisfies
\begin{equation}\label{eq:common-embedding}
 (1-\varepsilon)F(z)\le U_{G,q}(z)\le(1+\varepsilon)F(z)
 \qquad(z\in\R^v),
\end{equation}
then
\begin{equation}\label{eq:common-necessary}
 \bigl((1+\varepsilon)D\bigr)^{-1/p}
 \le4\varepsilon+\frac{2(1+\varepsilon)D}{R},
\end{equation}
and
\begin{equation}\label{eq:common-degree-bound}
 D\ge\frac1{1+\varepsilon}
       \min\{R\varepsilon,(6\varepsilon)^{-p}\}.
\end{equation}
\end{theorem}

\begin{proof}
The singleton inequalities \eqref{eq:weighted-degrees} specialize to
\begin{equation}\label{eq:unweighted-degrees}
 (1-\varepsilon)R\le\frac{\deg_G(u)}q
 \le(1+\varepsilon)R.
\end{equation}
Thus the support degree $k=\deg_G(u)$ obeys
$1\le k\le(1+\varepsilon)D$. Theorem~\ref{thm:weighted} gives
\[
 \bigl((1+\varepsilon)D\bigr)^{-1/p}
 \le k^{-1/p}
 \le4\varepsilon+\frac{2k}{R}
 \le4\varepsilon+\frac{2(1+\varepsilon)D}{R},
\]
which proves \eqref{eq:common-necessary}. Alternatively, combining
\eqref{eq:support-degree} with $k\le(1+\varepsilon)D$ directly gives
\eqref{eq:common-degree-bound}.
\end{proof}

\begin{corollary}[Even one successful realization is costly]
\label{cor:common-accuracy}
Fix $p>2$. Under \eqref{eq:large-v}, every common probability $q$ that
admits even one subgraph satisfying \eqref{eq:common-embedding} obeys
\begin{align}
 nq
 &\ge\frac{v}{2(1+\varepsilon)}(6\varepsilon)^{-p}
 \ge\frac{6^{-p}}3\,v\varepsilon^{-p},
 \label{eq:common-v-lower}\\
 nq
 &\ge\gamma_p(\Sp+d)\varepsilon^{-p},
 \qquad
 \gamma_p=\frac{6^{-p}}{3(2^{p-2}+1)}.
 \label{eq:common-sensitivity-lower}
\end{align}
If $nq$ violates the first lower bound in
\eqref{eq:common-v-lower}, no subgraph can succeed; the success
probability of the corresponding sampler is zero.
\end{corollary}

\begin{proof}
Use \eqref{eq:common-degree-bound}, \eqref{eq:large-v},
$nq=vD/2$, $1+\varepsilon<3/2$, and
\eqref{eq:sensitivity-comparison}. The necessary inequality constrains
the fixed parameter $q$ for each individual successful graph.
\end{proof}

Unlike \eqref{eq:weighted-expectation}, this bound has no factor
$1-\delta$. The common rescaling links every successful graph's degrees
to the prescribed $q$, so the existence of a single successful graph
already forces the expected-size bound. No realized edge count is
substituted for its expectation.

\section{Refutation of RA-06}\label{sec:refutation}

By Lemma~\ref{lem:sensitivity}, every row of $A_v$ has the same
sensitivity. Therefore each choice of $\alpha>0$ in
\eqref{eq:ra06-probabilities} produces one common probability
\begin{equation}\label{eq:equal-probability}
 q=\min\{1,1/n+s_v/\alpha\}>0.
\end{equation}
Its rescaled row energy is exactly $U_{G,q}$. The results of
Section~\ref{sec:common} cover every $q\in(0,1]$, including the
saturated choice $q=1$, and hence every probability allowed by
\eqref{eq:equal-probability}.

\begin{theorem}[Negative answer to RA-06]\label{thm:ra06}
For every fixed real $p>2$, there are no constants $C_p,c_p>0$
satisfying the target assertion of Section~\ref{sec:introduction}.
\end{theorem}

\begin{proof}
Suppose such constants exist. For an integer $b\ge3$, choose
\begin{equation}\label{eq:counterexample-sequence}
 \varepsilon=b^{-1},\qquad\delta=\tfrac14,
 \qquad v=\lceil b^{p+2}\rceil,\qquad A=A_v.
\end{equation}
For all sufficiently large $b$, one has
$(v-1)\varepsilon\ge(6\varepsilon)^{-p}$, so
Corollary~\ref{cor:common-accuracy} applies. Positive success probability
requires
\begin{equation}\label{eq:contradiction-lower}
 nq\ge\frac{6^{-p}}3\,v b^p.
\end{equation}
On the other hand, \eqref{eq:ra06-budget} and
Lemma~\ref{lem:sensitivity} require
\begin{equation}\label{eq:contradiction-upper-raw}
 nq\le C_p(2^{p-2}+1)\,vb^2
       \log^{c_p}\!\left(8b\binom v2(v-1)\right).
\end{equation}
Since $v\le2b^{p+2}$ for $b\ge3$,
\[
 8b\binom v2(v-1)\le4bv^3\le32b^{3p+7}.
\]
Thus there is a constant $K_{p,C_p,c_p}$, independent of $b$, such that
\begin{equation}\label{eq:contradiction-upper}
 nq\le K_{p,C_p,c_p}\,vb^2(\log b)^{c_p}.
\end{equation}
For every fixed $p>2$ and fixed $c_p>0$,
\[
 \frac{b^{p-2}}{(\log b)^{c_p}}\longrightarrow\infty.
\]
Consequently, \eqref{eq:contradiction-lower} and
\eqref{eq:contradiction-upper} are incompatible for large $b$.

Every $\alpha$ satisfying the proposed expected-size budget then has
success probability zero, whereas the assertion requires success
probability at least $3/4$. This proves the contradiction.
\end{proof}

The matrix may depend on $\varepsilon$: the target quantifies over
every matrix and every admissible accuracy. The sequence
\eqref{eq:counterexample-sequence} also keeps $n$, $d$, and
$\varepsilon^{-1}$ polynomially related, which is why no fixed
polylogarithmic factor can hide the accuracy gap.

\begin{corollary}[The same budget fails for arbitrary reweighted subsets]
\label{cor:general-budget}
Fix $p>2$ and a success probability $\rho\in(0,1]$. No constants
$C,c>0$ guarantee, for every full-column-rank matrix and every
$0<\varepsilon<1/2$, a randomized nonnegative weighted row subset
satisfying \eqref{eq:matrix-embedding} with probability at least $\rho$
and expected support size at most
\[
 C\varepsilon^{-2}(\Sp+d)
 \log^c\!\left(\frac{2nd}{\varepsilon}\right).
\]
This remains false when selection is dependent and the weights are
chosen adaptively.
\end{corollary}

\begin{proof}
For the sequence \eqref{eq:counterexample-sequence},
\eqref{eq:weighted-expectation} gives
$\E M\ge\rho\,6^{-p}vb^p/2$. The proposed budget is at most
$K_{p,C,c}vb^2(\log b)^c$. The same polynomial-versus-logarithmic
comparison gives a contradiction.
\end{proof}

\section{Scope and relation to prior bounds}\label{sec:scope}

\paragraph{What the obstruction rules out.}
Theorem~\ref{thm:ra06} refutes the joint dependence
$\widetilde O_p(\varepsilon^{-2}(\Sp+d))$ in RA-06.
Corollary~\ref{cor:general-budget} goes beyond its particular sampling
rule: nonnegative row reweighting and dependent selection cannot restore
that budget. The lower bounds concern row-subset representations,
including repeated rows after their weights are combined. They do not
apply to unrestricted linear sketches whose output rows are arbitrary
linear combinations of the input rows.

\paragraph{The dimension regime matters.}
The simplified $\varepsilon^{-p}$ lower bound is asserted under
\eqref{eq:large-v}; the finite-$v$ statement is the minimum in
\eqref{eq:weighted-support-bound}. No $\varepsilon^{-p}$ lower bound
is claimed for all accuracies with $v$ fixed. At fixed accuracy the
construction gives a lower bound only linear in $\Sp+d$, so it does
not itself rule out nearly linear sensitivity dependence in that regime.
The deterministic argument holds for $p>1$, but its comparison with
quadratic inverse-accuracy dependence produces a contradiction only
when $p>2$.

\paragraph{Ordinary-sensitivity upper bounds.}
Woodruff--Yasuda's proved ordinary-sensitivity bound is
$\widetilde O_p(\varepsilon^{-2}\Sp^{2-2/p})$ for $p>2$
\cite[Theorem~1.5]{wy}. Its larger sensitivity exponent distinguishes
it from the conjectured bound rejected here. Our lower bound does not
contradict that theorem. The conjecture and the sampling model are due
to Woodruff--Yasuda; the deterministic support obstruction and its
consequences are the results proved in this note.

\paragraph{A weaker inverse-accuracy dependence.}
Under its stated convexity and sampling-envelope hypotheses,
Dai--Kosov--Murasko's Theorem~3.5 gives a sufficient sample-size
condition containing the factor
$(\varepsilon^{-1}\log\varepsilon^{-1})^p$ \cite{dkm}.
Thus that theorem does not assert the quadratic accuracy dependence
refuted here. We use it only for this comparison of accuracy exponents;
no transfer of its guarantee to the particular Bernoulli rule
\eqref{eq:ra06-probabilities} is asserted or needed.

\begin{thebibliography}{9}

\bibitem{wy}
David P. Woodruff and Taisuke Yasuda.
\emph{Sharper Bounds for $\ell_p$ Sensitivity Sampling}.
ICML 2023; arXiv:2306.00732v2 (2024).
Definition~1.1, Theorem~1.5, and Section~3.
\par
\url{https://arxiv.org/html/2306.00732v2}.

\bibitem{ra06}
\emph{Open Problems in Numerical Linear Algebra}.
RA-06: Sensitivity-dependent row sampling for $\ell_p$ embeddings
when $p$ exceeds two.
Statement accessed 30 September 2026.
\par
\url{https://github.com/ajt60gaibb/OpenProblemsInNLA/blob/main/randomized-and-low-rank-approximation/RA-06/README.md}.

\bibitem{dkm}
Feng Dai, Egor Kosov, and Noel Murasko.
\emph{Empirical Approximation of $L_p$ Norms}.
arXiv:2606.00347v1 (2026), Theorem~3.5.
\par
\url{https://arxiv.org/html/2606.00347v1}.

\end{thebibliography}
\end{document}
